\section{What DEMETER Reveals about Train-and-Grow Optimization}

\begin{frame}[t]{What DEMETER Reveals about Train-and-Grow Optimization}
  \vspace{-0.08cm}
  \begin{columns}[T,onlytextwidth]
    \begin{column}{0.37\textwidth}
      {\fontsize{7.4}{9}\selectfont\bfseries\color{EnstaBlue}
       WHY ALIGNMENT BREAKS}

      \vspace{0.15cm}
      \begin{itemize}
        \setlength\itemsep{0.12cm}
        \fontsize{8.2}{10.2}\selectfont

        \item \textbf{Growth is function-directed.}
        New capacity is selected to align a local functional-gradient
        projection.

        \item \textbf{Training is parameter-directed.}
        SGD moves in parameter space; its Jacobian-induced effect generally
        follows a different function-space direction.

        \item \textbf{The directions can conflict.}
        Interleaving them can redirect progress; overtraining between growth
        events may favor local minima \textcolor{MutedText}{(Douka et al., 2025)}.
      \end{itemize}

      \vspace{0.08cm}
      \begin{beamercolorbox}[rounded=true,sep=0.17cm,wd=\linewidth]{demeter takeaway}
        \fontsize{7.6}{9.5}\selectfont
        \textbf{Key observation.} Alignment is episodic, not sustained.
        Cosine peaks mark growth epochs; between them, parametric training
        remains weakly aligned.
      \end{beamercolorbox}
    \end{column}

    \begin{column}{0.60\textwidth}
      \centering
      \includegraphics[width=\linewidth]{demeter_cosine_similarity.png}

      \vspace{0.05cm}
      \parbox{0.96\linewidth}{%
        \centering
        \fontsize{6.5}{8.0}\selectfont\color{MutedText}
        \textbf{Figure 4:} MultiNIST cosine similarity between the functional
        gradient and the zeroth-order contribution during DEMETER training and
        growth.%
      }
    \end{column}
  \end{columns}
\end{frame}
